\chapter{Reference Model}
\label{app:reference-model}

This appendix gives the complete model of Sections~\ref{sec:generic-model}
and~\ref{sec:harm-trees}. Its source is the file
\texttt{04\_Study/4.3\_Analytical-Framework/reference\_model\_v2.yaml}; the tables are generated from it.
\highlight[orange!25]{Version~2 differs from version~1 in two
places: K10 is an OR of two sufficient combinations, each naming its actor, instead of an OR+;
and deletion (IM12) and export (IM14) are separate components, each present only if no route
offers the function.}\todo{Added 4 November 2026 with reference model v2 (D12); version~1 is
kept unchanged in the same folder.}

\section{Components}
Each component is a feature that can be checked on an app. The colour of its ID shows its origin:
\ocore{core functionality}, \octx{digital context}, \oimp{implementation choice} or \oext{an event
outside the app}.

{\footnotesize
\begin{longtable}{@{}lp{4.1cm}p{1.9cm}p{6.1cm}@{}}
\caption[Components of the generic health tracker]{Components of the generic health tracker.}
\label{tab:reference-components}\\
\toprule
\textbf{ID} & \textbf{Component} & \textbf{Origin} & \textbf{Feature to check} \\
\midrule
\endfirsthead
\toprule
\textbf{ID} & \textbf{Component} & \textbf{Origin} & \textbf{Feature to check} \\
\midrule
\endhead
% Generated by scripts/instrument-test/v001/build_tables.py; do not edit.
\ocore{CO1} & Health records are logged or sensed & core & The app records the user's own health or body data (eligibility rule E1). \\
\ocore{CO2} & Insights are derived from the records & core & The app computes predictions, scores or insights from the records. \\
\ocore{CO3} & A long history is kept & core & Records accumulate over time, on the phone or on a server. \\
\octx{CX1} & Another person can use the phone & context & The records can be read in the app, so anyone who can open the phone can read them. \\
\octx{CX2} & Device and network data reach servers & context & IP address and device data reach a server when the app is online. \\
\octx{CX3} & Device or advertising identifier exists & context & The app can read or receive a persistent device, push or advertising identifier. \\
\oimp{IM1} & Account links records to a person & implementation & An account holds a name, e-mail address or similar next to the records. \\
\oimp{IM2} & Copy on the provider's servers & implementation & Records are stored or synchronised on servers. \\
\oimp{IM3} & Analytics or monitoring service & implementation & A third-party analytics or monitoring service receives usage data or identifiers. \\
\oimp{IM4} & Advertising or attribution partner & implementation & An advertising or attribution partner receives identifiers or app events. \\
\oimp{IM5} & Processor receives health data & implementation & A messaging, support, AI or other processor receives tracked health data. \\
\oimp{IM6} & De-identified data leaves the provider & implementation & De-identified or pseudonymised records go to research or other partners. \\
\oimp{IM7} & Commercial use of tracked data & implementation & Tracked data is sold or drives offers, recommendations or partner products. \\
\oimp{IM8} & Precise location or routes recorded & implementation & The app records GPS positions or routes. \\
\oimp{IM9} & Sharing with other people & implementation & Records or profiles can be shared with other people (partner sharing, feeds, community). \\
\oimp{IM10} & Transfer outside the EU/EEA & implementation & Personal data can be transferred outside the EU or EEA. \\
\oimp{IM11} & Function tied to agreement & implementation & A function is available only if the user agrees to data sharing beyond what it needs. \\
\oimp{IM12} & No usable way to delete or export & implementation & The user has no working in-app way to delete or export the data. \\
\oimp{IM13} & Flow without opt-out & implementation & A data flow to a third party has no opt-out while the app is in use. \\
\oext{EX1} & An attacker gains access & external & Assumed possible for any stored data; the study does not test security. \\
\oext{EX2} & An authority demands or seizes data & external & Depends on the legal setting; uncertain unless the evidence shows it. \\
\oext{EX3} & An insurer, employer or other party obtains data & external & Depends on recipients and context; uncertain unless the evidence shows it. \\
\bottomrule
% the generated rows end with \bottomrule
\end{longtable}}

\section{Capabilities}
K1 to K8 keep the numbers of an earlier capability table of the literature, which the brief
review of Section~\ref{sec:lit-review} replaced; the numbers stay so that the codebook and the
stored labels keep their meaning. K3, the inference of keystrokes from the motion sensors of a
smartwatch, is excluded because it needs the analysis of a connected device
(Section~\ref{sec:scope}). K9, K10 and I1 to I3 are added. The last column gives the components each capability needs, with the operators
of Section~\ref{sec:harm-trees}.

{\footnotesize
\begin{longtable}{@{}lp{3.5cm}p{2.5cm}lp{4.7cm}@{}}
\caption[Capabilities of the reference model]{Capabilities of the reference model.
\texttt{K-ACC}: access to the data, added in codebook version 5.}
\label{tab:reference-capabilities}\\
\toprule
\textbf{ID} & \textbf{Capability} & \textbf{Actor} & \textbf{Type} & \textbf{Needs} \\
\midrule
\endfirsthead
\toprule
\textbf{ID} & \textbf{Capability} & \textbf{Actor} & \textbf{Type} & \textbf{Needs} \\
\midrule
\endhead
% Generated by scripts/instrument-test/v002/build_reference.py; do not edit.
K1 & Locate homes, routes and places & Other users, attackers & \texttt{K-INF} & \oimp{IM8} AND (\oimp{IM9} OR \oext{EX1}) \\
K2 & Infer personal attributes & Recipients of health data & \texttt{K-INF} & (\ocore{CO1} OR \ocore{CO2}) AND (\oimp{IM5} OR \oimp{IM6} OR \oimp{IM7}) \\
K4 & Re-identify or link individuals & Holders of de-identified data & \texttt{K-TRK} & \oimp{IM6} AND \ocore{CO3} \\
K5 & Profile users across services & Advertising and analytics partners & \texttt{K-TRK} & (\oimp{IM3} OR \oimp{IM4}) AND (\octx{CX2} OR \octx{CX3}) \\
K6 & Read stored or intercepted data & Attackers & \texttt{K-ACC} & \oext{EX1} AND (\oimp{IM1} OR \oimp{IM2} OR \oimp{IM3} OR \oimp{IM5} OR \oimp{IM10}) \\
K7 & Obtain records for proceedings & Law enforcement, courts & \texttt{K-ACC} & \oext{EX2} AND (\oimp{IM2} OR \octx{CX1}) \\
K8 & Use data against users in insurance, work or relationships & Insurers, employers & \texttt{K-ACC} & \oext{EX3} AND (\oimp{IM6} OR \oimp{IM7} OR \oimp{IM9}) \\
K9 & Read entries through the phone or a share & Partner, family, people nearby & \texttt{K-ACC} & \ocore{CO1} AND (\octx{CX1} OR \oimp{IM9}) \\
K10 & Sort users into classes & Provider (CO2 AND IM7); advertising partners (IM4 AND CX3) & \texttt{K-DIS} & (\ocore{CO2} AND \oimp{IM7}) OR (\oimp{IM4} AND \octx{CX3}) \\
I1 & Cannot use a function without agreeing & User & \texttt{K-INC} & \oimp{IM11} \\
I2 & Cannot delete or take out the data & User & \texttt{K-INC} & \oimp{IM12} OR \oimp{IM14} \\
I3 & Cannot stop a data flow & User & \texttt{K-INC} & \oimp{IM13} \\
\bottomrule
% the generated rows end with \bottomrule
\end{longtable}}

\section{Pathways}
Each row is one pathway: a capability that can lead to a harm. All pathways of one harm form its
tree. \emph{Basis}: \emph{documented} if a source states the link, \emph{argued} if the study argues it
from its sources.

{\footnotesize
\begin{longtable}{@{}lllp{5.2cm}lp{3.5cm}@{}}
\caption[Pathways of the reference harm trees]{Pathways of the reference harm trees.}
\label{tab:reference-pathways}\\
\toprule
\textbf{Harm} & \textbf{ID} & \textbf{Cap.} & \textbf{Possible consequence} & \textbf{Basis} & \textbf{Sources} \\
\midrule
\endfirsthead
\toprule
\textbf{Harm} & \textbf{ID} & \textbf{Cap.} & \textbf{Possible consequence} & \textbf{Basis} & \textbf{Sources} \\
\midrule
\endhead
% Generated by scripts/instrument-test/v001/build_tables.py; do not edit.
\texttt{H-ECO} & P01 & K6 & Fraud or extortion after exposed data & documented & \parencite[App.~A]{knabenhans2026harms}\parencite[p.~5]{gabriele2020understanding} \\
 & P02 & K8 & Higher premiums, lost cover or lost job & documented & \parencite[p.~5]{gabriele2020understanding}\parencite[p.~3]{hudig2025intimate} \\
\addlinespace[2pt]
\texttt{H-PHY} & P03 & K1 & Stalking or burglary from known places & documented & \parencite[p.~14]{niksirat2024wearable}\parencite[App.~A]{knabenhans2026harms} \\
 & P04 & K9 & Violence or stalking by a partner & documented & \parencite[p.~14]{niksirat2024wearable}\parencite[p.~2]{malki2024exploring} \\
\addlinespace[2pt]
\texttt{H-PSY} & P05 & K5 & Fear of being watched & documented & \parencite[App.~A]{knabenhans2026harms} \\
 & P06 & K6 & Distress after exposure or identity theft & documented & \parencite[App.~A]{knabenhans2026harms} \\
 & P07 & K7 & Fear of legal consequences & documented & \parencite[pp.~1-2]{cao2024deleted} \\
 & P08 & K9 & Threats or harassment by someone close & documented & \parencite[p.~14]{niksirat2024wearable} \\
\addlinespace[2pt]
\texttt{H-AUT} & P09 & I1 & Coerced agreement & argued & \parencite[App.~A]{knabenhans2026harms}\parencite[p.~10]{malki2024exploring} \\
 & P10 & I2 & Locked into the service & argued & \parencite[p.~10]{malki2024exploring} \\
 & P11 & I3 & No choice over a data flow & argued & \parencite[p.~10]{malki2024exploring} \\
 & P12 & K7 & Restricted reproductive or health choices & documented & \parencite[p.~2]{hudig2025intimate}\parencite[p.~1]{cao2024deleted} \\
 & P13 & K10 & Manipulation through tailored offers & argued & \parencite[App.~A]{knabenhans2026harms} \\
\addlinespace[2pt]
\texttt{H-DSC} & P14 & K2 & Unequal treatment from inferred traits & documented & \parencite[p.~12]{niksirat2024wearable} \\
 & P15 & K8 & Unfair decisions on cover or hiring & documented & \parencite[p.~3]{hudig2025intimate}\parencite[p.~2]{malki2024exploring} \\
 & P16 & K10 & Profile-based prices or offers & documented & \parencite[App.~A]{knabenhans2026harms} \\
\addlinespace[2pt]
\texttt{H-REP} & P17 & K4 & Exposure of a re-identified person & argued & \parencite[Part~V.C]{citron2021harms} \\
 & P18 & K6 & Exposure of leaked records & argued & \parencite[Part~V.C]{citron2021harms} \\
 & P19 & K9 & Exposure to family or partners (\texttt{T-REL}) & documented & \parencite[p.~3]{hudig2025intimate} \\
\addlinespace[2pt]
\texttt{H-SOC} & P20 & K5 & Monitoring of groups of users & argued & \parencite[App.~A]{knabenhans2026harms} \\
 & P21 & K7 & Chilling effect of state access & documented & \parencite[App.~A]{knabenhans2026harms}\parencite[p.~2]{cao2024deleted} \\
\bottomrule
% the generated rows end with \bottomrule
\end{longtable}}

\section{Use Cases and Subtrees}
\label{sec:app-use-cases}
The harm trees of Section~\ref{sec:results-model} were built from the ten use cases below, in this
order (decision D23). Each card answers six questions: what the app \emph{records}, what it
\emph{derives}, who can reach the records on the \emph{phone}, what \emph{leaves} the phone and to
whom, what the app \emph{asks} of the user, and which event \emph{outside} the app a pathway
assumes. \emph{Basis}: documented if a source states it, argued if the study infers it. UC1 to UC7
build the trees; UC8 to UC10 are check cases. The table is generated from
\texttt{04\_Study/4.3\_Analytical-Framework/use\_cases\_v1.yaml}.

{\footnotesize
\begin{longtable}{@{}p{1.1cm}p{5.6cm}p{2.0cm}p{1.5cm}p{3.5cm}@{}}
\caption[Cards of the use cases]{Cards of the use cases.}
\label{tab:usecase-cards}\\
\toprule
\textbf{Question} & \textbf{Answer} & \textbf{Components} & \textbf{Basis} & \textbf{Evidence} \\
\midrule
\endfirsthead
\toprule
\textbf{Question} & \textbf{Answer} & \textbf{Components} & \textbf{Basis} & \textbf{Evidence} \\
\midrule
\endhead
% Generated by 04_Study/4.3_Analytical-Framework/scripts/harm-tree-usecases/v001/build_usecases.py; do not edit.
\textbf{UC1} & \textbf{Logging and predicting on the phone} \emph{(build)} & & & \\
record & Cycle, body and behaviour data entered by the user. & \ocore{CO1} & documented & \parencite[p.~3]{hudig2025intimate}; \parencite[p.~11]{shipp2020period} \\
derive & Predictions of cycle regularity and fertility windows. & \ocore{CO2} & documented & \parencite[p.~3]{hudig2025intimate} \\
record & Predictions rest on a history that grows with every entry. & \ocore{CO3} & argued & \parencite[p.~3]{hudig2025intimate} \\
phone & Anyone who can open the phone can read the entries; abusive partners use trackers to stalk and threaten. & \octx{CX1} & argued & \parencite[p.~14]{niksirat2024wearable} \\
leaves & Nothing. & -- &  &  \\
asks & Nothing. & -- &  &  \\
outside & None. & -- &  &  \\
\midrule
\textbf{UC2} & \textbf{Account and a copy on the provider's servers} \emph{(build)} & & & \\
leaves & An account ties a name or e-mail address to the records. & \oimp{IM1} & documented & \parencite[p.~6]{shipp2020period} \\
leaves & Records go to the provider's servers (14 of 30 menstrual apps; 15 of 19 policies name company servers). & \oimp{IM2} & documented & \parencite[p.~11]{shipp2020period}; \parencite[p.~8]{malki2024exploring} \\
leaves & Data can be transferred and kept outside the user's country; endpoints lie in several countries. & \oimp{IM10} & documented & \parencite[p.~9]{malki2024exploring}; \parencite[p.~1]{hudig2025intimate} \\
outside & Users fear that the provider's platform is hacked; the study assumes this for any stored data. & \oext{EX1} & argued & \parencite[p.~14]{niksirat2024wearable} \\
outside & 13 of 19 policies allow access by law enforcement on request or subpoena; users fear access by authorities. & \oext{EX2} & documented & \parencite[p.~8]{malki2024exploring}; \parencite[p.~1]{cao2024deleted} \\
\midrule
\textbf{UC3} & \textbf{Usage analytics and advertising} \emph{(build)} & & & \\
leaves & Usage data goes to an analytics or monitoring service. & \oimp{IM3} & documented & \parencite[p.~1]{huckvale2019assessment}; \parencite[p.~7]{malki2024exploring} \\
leaves & Identifiers or app events go to an advertising or attribution partner (29 of 36 apps sent data to Google and Facebook). & \oimp{IM4} & documented & \parencite[p.~1]{huckvale2019assessment}; \parencite[p.~9]{shipp2020period} \\
phone & IP address and device data reach a server whenever the app is online. & \octx{CX2} & documented & \parencite[p.~7]{malki2024exploring} \\
phone & A persistent device, push or advertising identifier exists and is read. & \octx{CX3} & documented & \parencite[p.~7]{malki2024exploring} \\
\midrule
\textbf{UC4} & \textbf{Sharing with other people} \emph{(build)} & & & \\
leaves & Data is shared with partners (for example in pregnancy plans), friends, family and strangers; caregivers monitor children. & \oimp{IM9} & documented & \parencite[p.~3]{hudig2025intimate}; \parencite[p.~8]{niksirat2024wearable} \\
outside & Insurers and employers are named as parties that could obtain the data. & \oext{EX3} & documented & \parencite[p.~3]{hudig2025intimate}; \parencite[p.~12]{niksirat2024wearable} \\
\midrule
\textbf{UC5} & \textbf{Recording routes and places} \emph{(build)} & & & \\
record & GPS coordinates and routes; locations can be inferred, and users rate location the most sensitive data (burglary, bike theft). & \oimp{IM8} & documented & \parencite[p.~12]{niksirat2024wearable}; \parencite[p.~14]{niksirat2024wearable}; \parencite[p.~1]{gabriele2020understanding} \\
\midrule
\textbf{UC6} & \textbf{Processors, research partners and commercial use} \emph{(build)} & & & \\
leaves & Support, messaging, AI or hosting processors receive tracked health data; 17 of 19 policies name third-party servers. & \oimp{IM5} & documented & \parencite[p.~3]{malki2024exploring}; \parencite[p.~8]{malki2024exploring} \\
leaves & De-identified records go to research and development or to partners; users can be offered opt-ins for research. & \oimp{IM6} & documented & \parencite[p.~8]{malki2024exploring}; \parencite[p.~5]{hudig2025intimate} \\
leaves & Data is shared with advertisers or sold; profiles drive tailored offers. & \oimp{IM7} & documented & \parencite[p.~2]{hudig2025intimate}; \parencite[p.~9]{shipp2020period} \\
\midrule
\textbf{UC7} & \textbf{Terms of use: consent, deletion, export and opt-out} \emph{(build)} & & & \\
asks & Refusing the terms often meant that the user could not use the service; users face a forced choice. & \oimp{IM11} & documented & \parencite[p.~9]{malki2024exploring}; \parencite[p.~10]{malki2024exploring}; \parencite[p.~14]{niksirat2024wearable} \\
asks & Only a third of the apps offered deletion in the app; one did not tell the server to delete backups. & \oimp{IM12} & documented & \parencite[p.~9]{shipp2020period}; \parencite[p.~3]{malki2024exploring} \\
asks & Portability of the data is rare. & \oimp{IM14} & documented & \parencite[p.~9]{shipp2020period}; \parencite[p.~2]{malki2024exploring} \\
asks & Users lack meaningful ways to opt out of practices while using the app. & \oimp{IM13} & documented & \parencite[p.~3]{malki2024exploring}; \parencite[p.~14]{niksirat2024wearable} \\
\midrule
\textbf{UC8} & \textbf{Diet and weight logging} \emph{(check)} & & & \\
record & Nutrition, weight and activity are tracked; eating and drinking can be inferred. & \ocore{CO1}, \ocore{CO3} & documented & \parencite[p.~1]{gabriele2020understanding}; \parencite[p.~12]{niksirat2024wearable} \\
derive & Progress and weight-loss insights. & \ocore{CO2} & documented & \parencite[p.~6]{niksirat2024wearable} \\
leaves & Account, server copy, analytics, advertising, commercial use and sharing in a community: the flows of UC2 to UC6. & \oimp{IM1}, \oimp{IM2}, \oimp{IM3}, \oimp{IM4}, \oimp{IM7}, \oimp{IM9} & argued & \parencite[p.~8]{niksirat2024wearable}; \parencite[p.~1]{huckvale2019assessment} \\
phone & Same phone and network conditions as UC1 and UC3. & \octx{CX1}, \octx{CX2}, \octx{CX3} & argued & \parencite[p.~7]{malki2024exploring} \\
\midrule
\textbf{UC9} & \textbf{Mood and cessation tracking} \emph{(check)} & & & \\
record & Mood, symptoms and cessation data. & \ocore{CO1}, \ocore{CO3} & documented & \parencite[p.~1]{huckvale2019assessment} \\
derive & Scores or feedback computed from the entries. & \ocore{CO2} & argued & \parencite[p.~1]{huckvale2019assessment} \\
leaves & Data goes to the developer's servers and to analytics and advertising services (transmission to a third party in 33 of 36 apps); chat transcripts of mental-health apps were mined for marketing. & \oimp{IM2}, \oimp{IM3}, \oimp{IM4}, \oimp{IM5}, \oimp{IM7} & documented & \parencite[p.~1]{huckvale2019assessment}; \parencite[p.~3]{malki2024exploring} \\
\midrule
\textbf{UC10} & \textbf{Sleep and heart rate synced from a wearable} \emph{(check)} & & & \\
record & Sensor records of sleep, heart rate and steps arrive in the app. & \ocore{CO1}, \ocore{CO3} & documented & \parencite[p.~6]{niksirat2024wearable}; \parencite[p.~1]{gabriele2020understanding} \\
leaves & The app side follows UC2 to UC4. & \oimp{IM2}, \oimp{IM3}, \oimp{IM4}, \oimp{IM9} & argued & \parencite[p.~8]{niksirat2024wearable} \\
phone & Same as UC1. & \octx{CX1} & argued & \parencite[p.~14]{niksirat2024wearable} \\
excluded & Inference of keystrokes or PINs from the motion sensors of the watch (K3): needs the analysis of the connected device, which the study excludes. &  &  &  \\
\bottomrule
% the generated rows end with \bottomrule
\end{longtable}}

\subsection*{Subtrees of UC4 to UC7}
Drawn as Figure~\ref{fig:usecase-uc1}. A capability that gains only a route and no pathway is not
drawn.

\begin{figure}[H]
\centering
% Generated by 04_Study/4.3_Analytical-Framework/scripts/harm-tree-usecases/v001/build_usecases.py; do not edit.
\begin{tikzpicture}[font=\scriptsize,x=1cm,y=1cm,hxharm/.style={draw=black!50,fill=black!6,rounded corners=2pt,align=center,text width=2.1cm,minimum height=8mm,inner sep=2pt},hxcap/.style={draw=accent,fill=accent!10,rounded corners=2pt,align=center,text width=2.3cm,minimum height=8mm,inner sep=2pt},hxgate/.style={draw=black!55,circle,inner sep=1.2pt,font=\tiny\bfseries,fill=white},hxleaf/.style={draw=#1,thick,rounded corners=1pt,align=center,fill=white,text width=2.05cm,minimum height=8mm,inner sep=2pt},hxed/.style={semithick,draw=black!55}]
  \node[hxharm] (HECO) at (1.10,0) {\textbf{Economic}};
  \node[hxharm] (HDSC) at (3.60,0) {\textbf{Discrimination}};
  \node[hxcap] (K8) at (2.35,-1.55) {\textbf{K8}\\Use data against users in insurance, work or relationships};
  \node[hxgate] (K8_n0) at (2.35,-3.10) {AND};
  \node[hxleaf=originext,fill=black!8] (K8_n1) at (1.18,-4.96) {\oext{EX3}\\An insurer, employer or other party obtains data};
  \node[hxleaf=originimp,fill=black!8] (K8_n2) at (3.53,-4.96) {\oimp{IM9}\\Sharing with other people};
  \draw[hxed] (K8.south) -- (K8_n0.north);
  \draw[hxed] (K8_n0.south) -- (K8_n1.north);
  \draw[hxed] (K8_n0.south) -- (K8_n2.north);
  \draw[hxed] (HECO.south) -- (K8.north) node[midway,fill=white,inner sep=0.8pt,font=\tiny]{P02};
  \draw[hxed] (HDSC.south) -- (K8.north) node[midway,fill=white,inner sep=0.8pt,font=\tiny]{P15};
\end{tikzpicture}

\caption[Subtree of UC4]{Subtree after UC4, sharing with other people.}
\label{fig:usecase-uc4}
\end{figure}

\begin{figure}[H]
\centering
% Generated by 04_Study/4.3_Analytical-Framework/scripts/harm-tree-usecases/v001/build_usecases.py; do not edit.
\begin{tikzpicture}[font=\scriptsize,x=1cm,y=1cm,hxharm/.style={draw=black!50,fill=black!6,rounded corners=2pt,align=center,text width=2.1cm,minimum height=8mm,inner sep=2pt},hxcap/.style={draw=accent,fill=accent!10,rounded corners=2pt,align=center,text width=2.3cm,minimum height=8mm,inner sep=2pt},hxgate/.style={draw=black!55,circle,inner sep=1.2pt,font=\tiny\bfseries,fill=white},hxleaf/.style={draw=#1,thick,rounded corners=1pt,align=center,fill=white,text width=2.05cm,minimum height=8mm,inner sep=2pt},hxed/.style={semithick,draw=black!55}]
  \node[hxharm] (HPHY) at (2.94,0) {\textbf{Physical}};
  \node[hxcap] (K1) at (2.94,-1.55) {\textbf{K1}\\Locate homes, routes and places};
  \node[hxgate] (K1_n0) at (2.94,-3.10) {AND};
  \node[hxleaf=originimp,fill=black!8] (K1_n1) at (1.18,-6.51) {\oimp{IM8}\\Precise location or routes recorded};
  \node[hxgate] (K1_n2) at (4.70,-4.65) {OR};
  \node[hxleaf=originimp] (K1_n3) at (3.53,-6.51) {\oimp{IM9}\\Sharing with other people};
  \node[hxleaf=originext] (K1_n4) at (5.88,-6.51) {\oext{EX1}\\An attacker gains access};
  \draw[hxed] (K1.south) -- (K1_n0.north);
  \draw[hxed] (K1_n0.south) -- (K1_n1.north);
  \draw[hxed] (K1_n2.south) -- (K1_n3.north);
  \draw[hxed] (K1_n2.south) -- (K1_n4.north);
  \draw[hxed] (K1_n0.south) -- (K1_n2.north);
  \draw[hxed] (HPHY.south) -- (K1.north) node[midway,fill=white,inner sep=0.8pt,font=\tiny]{P03};
\end{tikzpicture}

\caption[Subtree of UC5]{Subtree after UC5, recording routes and places.}
\label{fig:usecase-uc5}
\end{figure}

\begin{figure}[H]
\centering
% Generated by 04_Study/4.3_Analytical-Framework/scripts/harm-tree-usecases/v001/build_usecases.py; do not edit.
\scalebox{0.907}{%
\begin{tikzpicture}[font=\scriptsize,x=1cm,y=1cm,hxharm/.style={draw=black!50,fill=black!6,rounded corners=2pt,align=center,text width=2.1cm,minimum height=8mm,inner sep=2pt},hxcap/.style={draw=accent,fill=accent!10,rounded corners=2pt,align=center,text width=2.3cm,minimum height=8mm,inner sep=2pt},hxgate/.style={draw=black!55,circle,inner sep=1.2pt,font=\tiny\bfseries,fill=white},hxleaf/.style={draw=#1,thick,rounded corners=1pt,align=center,fill=white,text width=2.05cm,minimum height=8mm,inner sep=2pt},hxed/.style={semithick,draw=black!55}]
  \node[hxharm] (HDSC) at (5.29,0) {\textbf{Discrimination}};
  \node[hxharm] (HREP) at (14.40,0) {\textbf{Reputational}};
  \node[hxcap] (K2) at (5.29,-1.55) {\textbf{K2}\\Infer personal attributes};
  \node[hxcap] (K4) at (14.40,-1.55) {\textbf{K4}\\Re-identify or link individuals};
  \node[hxgate] (K2_n0) at (5.29,-3.10) {AND};
  \node[hxgate] (K2_n1) at (2.35,-4.65) {OR};
  \node[hxleaf=origincore] (K2_n2) at (1.18,-6.51) {\ocore{CO1}\\Health records are logged or sensed};
  \node[hxleaf=origincore] (K2_n3) at (3.53,-6.51) {\ocore{CO2}\\Insights are derived from the records};
  \node[hxgate] (K2_n4) at (8.23,-4.65) {OR};
  \node[hxleaf=originimp,fill=black!8] (K2_n5) at (5.88,-6.51) {\oimp{IM5}\\Processor receives health data};
  \node[hxleaf=originimp,fill=black!8] (K2_n6) at (8.23,-6.51) {\oimp{IM6}\\De-identified data leaves the provider};
  \node[hxleaf=originimp,fill=black!8] (K2_n7) at (10.58,-6.51) {\oimp{IM7}\\Commercial use of tracked data};
  \node[hxgate] (K4_n0) at (14.40,-3.10) {AND};
  \node[hxleaf=originimp,fill=black!8] (K4_n1) at (13.23,-6.51) {\oimp{IM6}\\De-identified data leaves the provider};
  \node[hxleaf=origincore] (K4_n2) at (15.58,-6.51) {\ocore{CO3}\\A long history is kept};
  \draw[hxed] (K2.south) -- (K2_n0.north);
  \draw[hxed] (K4.south) -- (K4_n0.north);
  \draw[hxed] (K2_n1.south) -- (K2_n2.north);
  \draw[hxed] (K2_n1.south) -- (K2_n3.north);
  \draw[hxed] (K2_n0.south) -- (K2_n1.north);
  \draw[hxed] (K2_n4.south) -- (K2_n5.north);
  \draw[hxed] (K2_n4.south) -- (K2_n6.north);
  \draw[hxed] (K2_n4.south) -- (K2_n7.north);
  \draw[hxed] (K2_n0.south) -- (K2_n4.north);
  \draw[hxed] (K4_n0.south) -- (K4_n1.north);
  \draw[hxed] (K4_n0.south) -- (K4_n2.north);
  \draw[hxed] (HDSC.south) -- (K2.north) node[midway,fill=white,inner sep=0.8pt,font=\tiny]{P14};
  \draw[hxed] (HREP.south) -- (K4.north) node[midway,fill=white,inner sep=0.8pt,font=\tiny]{P17};
\end{tikzpicture}}

\caption[Subtree of UC6]{Subtree after UC6, processors, research partners and commercial use.}
\label{fig:usecase-uc6}
\end{figure}

\begin{figure}[H]
\centering
% Generated by scripts/harm-tree-usecases/v001/build_usecases.py; do not edit.
\begin{tikzpicture}[font=\scriptsize,x=1cm,y=1cm,hxharm/.style={draw=black!50,fill=black!6,rounded corners=2pt,align=center,text width=2.1cm,minimum height=8mm,inner sep=2pt},hxcap/.style={draw=accent,fill=accent!10,rounded corners=2pt,align=center,text width=2.3cm,minimum height=8mm,inner sep=2pt},hxgate/.style={draw=black!55,circle,inner sep=1.2pt,font=\tiny\bfseries,fill=white},hxleaf/.style={draw=#1,thick,rounded corners=1pt,align=center,fill=white,text width=2.05cm,minimum height=8mm,inner sep=2pt},hxed/.style={semithick,draw=black!55}]
  \node[hxharm] (HAUT) at (5.00,0) {\textbf{Autonomy}};
  \node[hxcap] (I1) at (1.18,-1.55) {\textbf{I1}\\Cannot use a function without agreeing};
  \node[hxcap] (I2) at (5.00,-1.55) {\textbf{I2}\\Cannot delete or take out the data};
  \node[hxcap] (I3) at (8.82,-1.55) {\textbf{I3}\\Cannot stop a data flow};
  \node[hxleaf=originimp,fill=black!8] (I1_n0) at (1.18,-4.96) {\oimp{IM11}\\Function tied to agreement};
  \node[hxgate] (I2_n0) at (5.00,-3.10) {OR};
  \node[hxleaf=originimp,fill=black!8] (I2_n1) at (3.83,-4.96) {\oimp{IM12}\\No usable way to delete the data};
  \node[hxleaf=originimp,fill=black!8] (I2_n2) at (6.17,-4.96) {\oimp{IM14}\\No usable way to export the data};
  \node[hxleaf=originimp,fill=black!8] (I3_n0) at (8.82,-4.96) {\oimp{IM13}\\Flow without opt-out};
  \draw[hxed] (I1.south) -- (I1_n0.north);
  \draw[hxed] (I2.south) -- (I2_n0.north);
  \draw[hxed] (I3.south) -- (I3_n0.north);
  \draw[hxed] (I2_n0.south) -- (I2_n1.north);
  \draw[hxed] (I2_n0.south) -- (I2_n2.north);
  \draw[hxed] (HAUT.south) -- (I1.north) node[midway,fill=white,inner sep=0.8pt,font=\tiny]{P09};
  \draw[hxed] (HAUT.south) -- (I2.north) node[midway,fill=white,inner sep=0.8pt,font=\tiny]{P10};
  \draw[hxed] (HAUT.south) -- (I3.north) node[midway,fill=white,inner sep=0.8pt,font=\tiny]{P11};
\end{tikzpicture}

\caption[Subtree of UC7]{Subtree after UC7, terms of use.}
\label{fig:usecase-uc7}
\end{figure}
